\chapter*{Abstract}
\addcontentsline{toc}{chapter}{Abstract}

Intermodal container terminals coordinate the transfer of containers between
rail, road and storage. Jointly scheduling crane operations, truck parking and
train loading is a combinatorial problem that operations research has
traditionally decomposed into narrowly defined sub-problems. This thesis
investigates whether a single deep reinforcement learning (DRL) agent can
instead learn to coordinate these tightly coupled tasks jointly. The objective
is not to outperform existing terminal operations, but to establish whether
such joint coordination is feasible.

The agent represents the terminal as a multi-channel spatial tensor processed
by a 3D convolutional neural network. Actions are expressed in a factored
source--destination space with dynamic masking to enforce physical validity.
Learning is guided by a curriculum of 33 tutorial scenarios progressing from
primitives to full logistics workflows. All experiments are grounded in
operational data from an intermodal terminal in N\"urnberg, Germany.
Candidate architectures are first screened on the tutorial curriculum before
deployment in controlled scalability experiments. The screening reveals that
compact convolutional backbones outperform larger alternatives and that
rankings reverse under operational time constraints; multi-seed validation
confirms the robustness of the final selection. In scalability experiments,
agents continue learning online as workload size increases. Task completion remains $\geq 99\%$ for workloads up to
200~containers (factor 10 scale-up from tutorial) while a priority-based heuristic 
baseline achieves only 52.4\%. A mixed-operation stress test further demonstrates 
simultaneous coordination of all eight container move types with $\geq 99.7\%$ completion, 
indicating that the agents acquire meaningful spatial reasoning over terminal operations.

These results provide a qualified \emph{yes}: a single DRL agent can learn the
joint coordination problem when supported by a structured spatial
representation, a factored action space and a progressive curriculum.
However, the learned policy behaves primarily as a \emph{reactive dispatcher},
making effective local decisions but lacking the ability to maintain
multi-step plans or perform proactive yard rearrangement. Rather than claiming
operational optimality, this work establishes the feasibility of DRL-based
joint coordination and lays groundwork for future research on more strategic,
planning-capable agents.
